\section{Exact zero-flow projection spectrum}
\label{supp:phasespectrum}
At unit total density, zero force and zero velocity, the projection
recoloring map gives outgoing phase populations
\begin{equation}
 f_i^{0,*}=w_i[\varphi+\beta\varphi(1-\varphi)\bm e_i\cdot\bm n_h].
\end{equation}
Their sum after streaming defines a nonlinear scalar map $P(\varphi)$.
Neutral ghost reflection and conservative phase transfer reproduce its
reflected scalar extension at the wall. This is the projection model,
not the older cosine normalization.

Let $p_j$ be the independently relaxed periodic planar profile,
$g_j=(p_{j+1}-p_{j-1})/2$, and $n_j=\operatorname{sgn}(g_j)$ where
$2|g_j|>10^{-9}$, zero otherwise. The factor two matters because production
applies its cutoff to the color gradient of $\chi=2\varphi-1$.
Define shifts $(Lv)_j=v_{j-1}$ and $(Rv)_j=v_{j+1}$, and diagonal operators
\begin{equation}
 E=\operatorname{diag}[(1-2p_j)n_j],\quad
 C=\operatorname{diag}\left[\frac{p_j(1-p_j)}{|g_j|}\right]_{active},
 \quad S=\frac23 I+\frac16(L+R).
\end{equation}
Inactive entries of $C$ are zero. For a transverse Fourier perturbation
$\delta\varphi_j\exp(iqy)$, differentiation of the gradient gives
$\delta g_y=i\sin(q)S\delta\varphi$ and
$\delta n_y=\delta g_y/|g_j|$ on the active mask. Variation of $n_x$
vanishes to first order around the axial normal. With $a=1/9$, $b_d=1/36$
and $c=\cos q$, summing all nine links yields
\begin{align}\label{eq:exactsymbol}
 T(q)={}&(4/9+2ac)I+(a+2b_dc)(L+R)\nonumber\\
 &+\beta(a+2b_dc)(LE-RE)\nonumber\\
 &+2\beta\sin^2(q)[aI+b_d(L+R)]CS.
\end{align}
The last term is the unit-normal variation; omitting it would give a
different leading shape dynamics. The profile and active mask are common
to the scalar matrix and production calculation. The derivative of a
hard cutoff is considered only within this fixed active-mask neighborhood.
No differentiability claim at the cutoff transition itself is made.

The neutral reflected transverse modes have $q=m\pi/H$. We select the
real positive eigenvalue whose eigenvector has the greatest normalized
overlap with the central derivative of $p$. This overlap is above 0.9979
in the measured cases. The physical rate is $-\log\mu/h^2$, not
$(1-\mu)/h^2$. Numerical differentiation of $P$ with amplitude $10^{-5}$
checks the full matrix action with maximum discrepancy
$\PhaseSymbolError$ over the three grids.

Full simulations start from outgoing populations built from
$p_j-\epsilon\cos[m\pi(y+1/2)/H]\,\partial_xp_j$, with physical
$\epsilon=\pm0.01$, and $\pm0.02$ at H16. They use zero surface tension,
zero body force and no hydraulic slip. A phase integral over one interface,
projected on the transverse cosine, measures its displacement. The decay
rate is diagnosed by a log-linear fit over the latter half of each
2000-physical-time trajectory. This fit measures the simulation response;
it does not enter the matrix prediction.

\begin{center}\small
\begin{tabular}{rrrr}\toprule
$H$ & $m$ & Predicted physical rate & Maximum relative difference \\\midrule
16 & 1 & $6.09541216\times10^{-5}$ & $7.392\times10^{-6}$ \\
32 & 1 & $1.55521488\times10^{-5}$ & $2.378\times10^{-6}$ \\
48 & 1 & $6.90314294\times10^{-6}$ & $2.454\times10^{-6}$ \\
16 & 2 & $9.86115881\times10^{-4}$ & $1.473\times10^{-6}$ \\
32 & 2 & $2.48320427\times10^{-4}$ & $2.932\times10^{-6}$ \\
\bottomrule\end{tabular}
\end{center}
Bounds round upward across both perturbation signs and tested amplitudes.
The spectral rate ratios support fourth-order shape relaxation in this
test. An arbitrary-interface modified equation, universal coefficient
or isotropy theorem is not inferred from five measured mode/grid combinations.
The known leading conservative Allen--Cahn mapping and projection/isotropy
distinction remain prior results \cite{subhedar2020}.

\subsection{Formal isolated-interface expansion}
The $h^2/24$ trend can be understood in an idealized limit, separately
from the exact finite-grid spectrum above. Consider one isolated ascending
interface in an unbounded continuum coordinate, smooth decaying
perturbations, projection weights and no active tail cutoff. Set
$\beta=bh$, $J=\varphi(1-\varphi)\bm n$, and
$p_0'=2bp_0(1-p_0)$. The D2Q9 moments give
\begin{align}
 \bm g_h&=\nabla\varphi+\frac{h^2}{6}\nabla\Delta\varphi+O(h^4),\\
 \frac{P(\varphi)-\varphi}{h^2}
 &=\frac16\Delta\varphi-\frac b3\nabla\cdot J
 +h^2\left[\frac1{72}\Delta^2\varphi
 -\frac b{18}\Delta\nabla\cdot J
 -\frac b3\nabla\cdot\{\varphi(1-\varphi)\bm n_2\}\right]+O(h^4),
\end{align}
where $\bm n_h=\bm n+h^2\bm n_2+O(h^4)$. The gradient normalization
is part of this expansion; $\bm n_2$ must not be discarded.

For $\delta\varphi=-\eta(y)p_0'(x)$,
$\delta n_y=-\eta'$ and $\delta n_{2,y}=-\eta'''/6$.
The $p_0'''$ terms in gradient normalization cancel between numerator
and denominator. On the leading profile,
$\nabla\cdot J=\Delta\varphi/(2b)$, so the first two fourth-order
terms combine to $-\Delta^2\varphi/72$.
Integration of their linear variation across the interface gives
$\eta''''/72$; the normal correction gives
$b\eta''''\int p_0(1-p_0)\,dx/18=\eta''''/36$.
The normal-profile correction does not introduce an integrated $k^2$
term at this order: its stationary equation gives
$p_h'-2bp_h(1-p_h)=h^2p_0'''/12+O(h^4)$, whose leading correction
integrates to zero. The leading linear phase operator annihilates
translation and conserves its integral, so the discrete-time logarithm
has no additional contribution to this projected rate at order $h^2$.
Since $\int\delta\varphi\,dx=-\eta$, the formal result is
\begin{equation}
 \partial_t\eta=-\frac{h^2}{24}\partial_y^4\eta+O(h^4),\qquad
 \Gamma=\frac{h^2 k_{phys}^4}{24}+O(h^4).
\end{equation}
This conditional calculation explains the reference trend without fitting
a coefficient. It is a formal smooth isolated-interface expansion, not
a uniform theorem with production cutoffs, two interacting interfaces,
nonneutral walls or arbitrary moving geometries. Those effects remain
in $T(q)$; its eigenvalues supply the reported predictions. A formal
coefficient is not substituted for the exact finite-grid comparison.
As a separate mathematical control, exact matrices at the analytic
zero-flux planar recurrence, diagnostic cutoff zero and fixed physical
$k=\pi/2$ have coefficient $\Gamma/(h^2k^4)$ equal to
0.03809364, 0.04072182 and 0.04142755 at H32/H64/H128.
Their relative differences from $1/24$ are $-8.575$\%, $-2.268$\%
and $-0.574$\%, with successive error orders 1.9190 and 1.9824.
The scalar stationary residual is below $4\times10^{-16}$.
These are additional matrix/expansion checks, not full production
trajectories with altered cutoffs. The production-spectrum command
also rejects a profile that has not relaxed sufficiently.

\section{Prescribed neutral slip and independent dynamic controls}
\subsection{Forward specular reflection and component transfer}
Let $\Delta f_A$ and $\Delta f_T$ denote differences between the appropriate
forward-neighbor outgoing specular link and the original local outgoing
bounce-back link, for phase A and total fluid. At a wall face, let
$\bar\varphi$ be the mean of its two outgoing node fractions. The new
incoming A population is
\begin{equation}
 f_A^{in}=f_A^{BB}+s\Delta f_A+
 (1-s)[\Delta f_A-\bar\varphi\Delta f_T],\qquad
 f_T^{in}=f_T^{BB}+s\Delta f_T.
\end{equation}
The B population is complementary. Each face has its opposite partner,
so component transfers sum to zero. Pure phases have zero relative transfer.
At rest with constant total density, $\Delta f_T=0$; the phase update is
the full scalar forward reflection independently of $s$.
Using a same-node specular value or leaving the relative transfer unscaled
does not give this closure. Earlier pilots with that combination are
retained and excluded from the final comparisons.

The admitted closure has $s=\lambda_l/(\lambda_l+1/2)$ and equal
$\tau=1$, legacy MRT, standard pressure weights, projection recoloring,
CSF and neutral flat ghost/flux walls. The three checked physical slip
lengths are 0.5, 1 and 2 at $H=8,16,32$, with both driving signs.
The independently calculated midpoint parabola supplies the hydraulic
reference. The nine CTest groups include eighteen pure profiles,
neutral fixed-point preservation, signed transport, rotated two-rank
exchange, one-/two-rank outgoing restart and invalid-configuration exits.
Three closed default settings retain bitwise fields and populations
against the preceding frozen kernel. Slip is a specified numerical
boundary control; a molecular slip parameter is not inferred.

\subsection{Staggered Stokes and finite-volume phase reference}
The independent hydraulic solver uses cell pressure, x-face axial
velocity and y-face transverse velocity. A corner streamfunction gives
$u=\delta_y\psi$ and $v=-\delta_x\psi$, with divergence cancelling
algebraically. Periodic x Fourier modes diagonalize the axial Laplacian;
the wall-normal biharmonic matrix solves
$\nu\Delta_h^2\psi=\operatorname{curl}_h\bm F$. Nonzero modes have
$\psi=0$ on both walls. Navier slip supplies wall vorticity through
the ghost factor $(1-2\lambda/h)/(1+2\lambda/h)$. The mean axial flow is
solved separately using ghost velocity ratio
$(2\lambda-h)/(2\lambda+h)$, so it is not lost in a periodic
streamfunction inversion.

For constant axial force, its exact sampled discrete profile is
\begin{equation}
 u_j=\frac{F}{2\nu}
 [y_j(16-y_j)+16\lambda+h^2/4].
\end{equation}
The positive $h^2/4$ difference from the LBM midpoint reference is retained.
All checked no-slip/Navier discrete profiles agree within $2\times10^{-12}$
relatively. Manufactured flow from
$\psi=y^2(16-y)^2\sin(2\pi x/64)$ has relative velocity errors
0.0248095, 0.00622349 and 0.00155717 on 16/32/64 grids, giving
orders 1.9951 and 1.9988. A discretely exact gradient force produces
maximum velocity below $5\times10^{-17}$; manufactured divergence is
also checked separately.

The independent scalar uses conservative centered face advection and
second-order Runge--Kutta stepping. Direct primitive-variable pilot
fluxes produced transverse distortion and incorrect flow direction;
reducing the time step did not resolve that failure. These failures
are retained and do not validate the leading PDE.
The accepted reference discretizes the continuously equivalent flux
using $d=\log[\varphi/(1-\varphi)]/(2b)$:
\begin{equation}
 \bm q=D\,2b\varphi(1-\varphi)(\nabla d-\bm n_d).
\end{equation}
On a face, $2b\varphi(1-\varphi)$ is represented by the divided
difference $\Delta\varphi/\Delta d$, with its analytic equal-value limit.
Tangential face gradients and unit normals use the distance field.
The clipping interval $[10^{-12},1-10^{-12}]$ regularizes reconstruction
only; the evolving fractions are not clipped. It is inactive throughout
the reported comparisons. Published signed-distance/profile-preservation
approaches motivate this representation \cite{hao2024}; their complete
algorithms and boundedness proofs are not claimed for our stencil.

For finite-volume runs, the analytic tanh coefficient is chosen to match
the same interpolated $W_{0.9}$ crossing width, independently of velocity.
The scalar keeps a flat tanh bridge stationary and translates a prescribed
profile at speed 0.001 over physical time 100 with relative centroid
errors $-1.381\times10^{-4}$ and $-3.651\times10^{-5}$ at H16/H32.
Two-phase zero-sigma runs recover the exact discrete pure hydraulics
within $2\times10^{-15}$ relatively. The dynamic reference has
$\nu=D=1/6$, $\sigma=.005$, physical force $10^{-7}$ and slip 1.
Its phase mass changes by less than $2\times10^{-13}$ relatively.
H32 negative-force half-step and positive-force full-step mean speeds
agree in magnitude within $\PhaseFVTimeSignError$. The last output-interval
changes are below $7\times10^{-6}$ relatively. These checks bound the
specified discretization; they do not establish full dynamic convergence.

At H32, an additional FV run starts from the independently generated
discrete projection profile, while an additional ablation starts from the
analytic logistic profile. The respective late mean-speed changes are
$-1.629\times10^{-6}$ and $3.249\times10^{-13}$ relatively; centroid
changes are $-1.783\times10^{-6}$ and $-6.352\times10^{-11}$.
This targeted cross-control distinguishes initialization from the changed
phase operator. It does not remove the dynamic continuum discrepancy.

The ablation adds $(P(\varphi)-\varphi)/h^2$ to the conservative
face-advection right-hand side in place of the finite-volume phase flux.
The Stokes solver and distance-based capillary force stay independent.
It is a semidiscrete shared-kinetics experiment, rather than a replay of
the full LBM or an independent PDE validation. Its phase mass changes
by less than $4\times10^{-12}$ relatively; the H32 half-step mean-flow
comparison is recorded separately. Mean and circular-centroid velocities
remain distinct observables in the data.

\subsection{Actual-force audit and retained evidence}
The production force audit uses the new frozen executable's own
gradient, curvature and cutoff on two adjacent outgoing checkpoints.
Its density/color match the physical fields within $3\times10^{-12}$.
Cell forces are averaged onto the appropriate MAC faces, converted by
$F_{phys}=F_l/h^3$, and combined with the independent uniform body force.
The actual-force and same-phase-field reference-force means are retained
for each state. These are hydraulic diagnostics on a supplied phase field;
their small residual does not independently verify its evolution.

\texttt{data/phase-kinetics.json} records 55 LBM executions, including
14 phase-wave trajectories, six default-compatibility runs and three
two-state hydraulic replays. It separately records seven independent
finite-volume runs and five shared-phase ablations. The collector
recomputes all diagnostics from raw fields, matches frozen reference
source hashes and checks inputs, geometry, executables and native
outgoing donor hashes. The preceding zero-slip kernel remains distinct
from the optional-slip kernel pinned by
\texttt{run-kernels/navier-slip-manifest.json}.

Original receipts retain diagnostic mistakes: early pure-run elapsed times
used step minus one although their start was zero; exploratory spectrum
files omitted the factor two in the active-gradient cutoff. The collector
recomputes both correctly. An early H48 ablation recorded its source hash
at completion after a docstring/scope edit; it is excluded and replaced
by a canonical rerun with launch-time source/dependency hashes. Earlier
ablation scopes are corrected explicitly in the scientific summary while
their original metadata are retained. Primitive-reference failures,
same-node slip pilots and unverified early reference source identities
remain outside the claimed successful comparisons.
